
\begin{exercise}[Infinite Monkey Theorem \Coffeecup\Coffeecup]
    The infinite monkey theorem states that an immortal monkey hitting keys at
    random on a typewriter keyboard for an infinite amount of time will
    eventually type any arbitrarily long text -- for example all the
    works of William Shakespeare -- with probability one.  It is also known as
    the ``law of truly large numbers''.
    Prove it!
    \begin{enumerate}[label={\emph{Level \arabic*:}},ref={Level \arabic*},leftmargin=4em,noitemsep]
        \item Assume the hit keys are independent and uniformly distributed.
        \item\label{it: level iid} Assume the hit keys are iid with a probability distribution that
        assigns each of the (finite) keys a positive probability.
        \item\label{it: level markov chain} Assume the hit keys are a stationary Markov chain whose transition
        probabilities from key \(x\) to key \(y\) are positive for all pairs
        \(x,y\) from the (finite) alphabet.
    \end{enumerate}
    \textbf{Hint:} \ref{it: level markov chain} is significantly harder than
    the previous levels. For it you already need to know Markov chains
    and modify the proof of the second Borel-Cantelli lemma for your purposes.
    You may want to come back to this later.
\end{exercise}

\begin{solution}
    Let \(\alphabet\) be a finite alphabet of keys and let \(X_n\in
    \alphabet\) be the letter typed at time \(n\). Let \(y\in \alphabet^m\)
    be the desired text of length \(m\). We essentially assume \(m=1\)
    without loss of generality with the definition of
    \[
        Y_n = (X_{nm}, X_{nm+1}, \ldots, X_{m(n+1)-1}) \in \alphabet^m,
    \]
    which is still a sequence of iid random variables in \ref{it: level iid}
    (or still a Markov chain in \ref{it: level markov chain}).  Let us consider \ref{it: level iid}
    first. Since
    \[\begin{aligned}
        \delta \coloneq \prob{Y_n = y}
        &= \prob{X_{nm} = y_1, X_{nm+1} = y_2, \ldots, X_{m(n+1)-1} = y_m}
        \\
        &= \prod_{i=1}^m \prob{X_{nm+i-1} = y_i} > 0,
    \end{aligned}
    \]
    Due to \(\sum_{n=1}^\infty \prob{Y_n = y} = \sum_{n=1}^\infty \delta = \infty\)
    we may apply the second Borel-Cantelli lemma \eqref{eq: borel cantelli 2} to get
    \[
        \prob[\big]{Y_n = y \text{ for infinitely many } n}
        = \prob[\big]{\limsup_{n\to\infty} \set{Y_n = y}}
        \overset{\text{\eqref{eq: borel cantelli 2}}}= 1.
    \]
    Observe that we have proven more than the desired claim: The monkey will
    not only type the desired text \(y\) at least once but infinitely many
    times with probability one!

    For the proof of \ref{it: level markov chain} we need to avoid independence.
    Since any transition probability is positive, we have
    \[
        \delta \coloneq \prob{Y_n = y \mid Y_{n-1} \neq y}
        = \sum_{\substack{x \in \alphabet^m \\ x \neq y}}\prob{Y_n = y \mid Y_{n-1} = x}
        > 0.
    \]
    We will now modify the proof of the second Borel-Cantelli lemma 
    in the place where independece was used. Specifically, we have
    \begin{equation}
        \label{eq: modified borel cantelli 2}
        \prob[\big]{(\limsup_{n\to\infty} \set{Y_n = y})^\complement}
        = \lim_{N\to \infty} \lim_{k\to \infty} \prob[\Big]{\bigcap_{n=N}^k \set{Y_n \neq y}}.
    \end{equation}
    with
    \[\begin{aligned}
        \prob[\Big]{\bigcap_{n=N}^k \set{Y_n \neq y}}
        &= \prod_{n=N}^k \prob{Y_n \neq y \mid Y_{n-1} \neq y,\ldots, Y_{N} \neq y}
        \\
        \overset{\text{Markov}}&= \prob{Y_n\neq y}\prod_{n=N+1}^k \prob{Y_n \neq y \mid Y_{n-1} \neq y}
        \\
        &\le (1-\delta)^{k-N-1}.
    \end{aligned}
    \]
    Thereby \eqref{eq: modified borel cantelli 2} is zero. Consequently,
    the probability that \(Y_n = y\) happens infinitely often is one.
\end{solution}
